% GENERATED FILE. Edit canonical lessons, then run npm run generate:books.
% canonical-source-hash: fnv1a64:45407c2d1f1293a0
% canonical-lessons: LA-C523-robustus, LA-C523-obesus, LA-C523-gracilis, LA-C523-sanitas, LA-C523-salubris

\chapter{Build and health}
\label{ch:la-build-and-health}
\hlchaptermodality{523}

\begin{chapteropening}
\textbf{By the end of this chapter:} \emph{I can describe someone's build and health.}

Complete the last lesson of chapter 523: I can describe someone's build and health.
\end{chapteropening}

\section[rōbustus, rōbusta, rōbustum]{rōbustus, rōbusta, rōbustum — sturdy, strong}

\label{lesson:LA-C523-robustus}



\begin{quote}
Before the new one: say the Latin for thin, then the Latin for leanness.
\end{quote}

\subsection*{You'll want to know: rōbustus, rōbusta, rōbustum}
\textbf{rōbustus, rōbusta, rōbustum} — \textquotedblleft{}sturdy, strong\textquotedblright{}.

\begin{quote}
Agricola rōbustus saccum gravem portat.
\end{quote}

\begin{quote}
\emph{The sturdy farmer carries a heavy sack.}
\end{quote}

From \emph{rōbur} — oak, hardwood: as strong as oak. English \emph{robust}.

\begin{cousinweb}
\textbf{rōbustus} is built on \emph{rōbur} \textquotedblleft{}oak, hard wood, strength\textquotedblright{} (Lewis \& Short; EtymDB): oaken, then sturdy. English \emph{robust} comes from it.
\end{cousinweb}

\begin{grammarlens}[title={What you've built}]
One more word for telling, explaining and describing at B1.
\end{grammarlens}

\subsection*{Your turn}
\begin{itemize}
\raggedright
  \item \emph{Say it:} \emph{rōbustus, rōbusta, rōbustum}
  \item \emph{Say it:} \emph{rōbustus, rōbusta, rōbustum}, once more
  \item \emph{Say it:} a sentence of your own with \emph{rōbustus, rōbusta, rōbustum}
  \item \emph{Recall:} say the Latin for thin, then the Latin for leanness, then say \emph{rōbustus, rōbusta, rōbustum} again
\end{itemize}

\begin{tcolorbox}[breakable,colback=teal!4,colframe=teal!35!black,title={Before you move on}]
What does rōbustus, rōbusta, rōbustum mean? (\textquotedblleft{}Sturdy, strong\textquotedblright{}.) Say it once more.
\end{tcolorbox}
\section[obēsus, obēsa, obēsum]{obēsus, obēsa, obēsum — fat, stout}

\label{lesson:LA-C523-obesus}



\begin{quote}
Before the new one: say the Latin for leanness, then the Latin for sturdy.
\end{quote}

\subsection*{You'll want to know: obēsus, obēsa, obēsum}
\textbf{obēsus, obēsa, obēsum} — \textquotedblleft{}fat, stout\textquotedblright{}.

\begin{quote}
Coquus obēsus multum edit.
\end{quote}

\begin{quote}
\emph{The fat cook eats a lot.}
\end{quote}

The past participle of \emph{obedō} — to eat away; here, well fed. English \emph{obese}.

\begin{cousinweb}
\textbf{obēsus} is the past participle of \emph{obedō}, built on \emph{ob-} and \emph{edō} \textquotedblleft{}to eat\textquotedblright{} (Lewis \& Short). English \emph{obese} comes from it.
\end{cousinweb}

\begin{grammarlens}[title={What you've built}]
One more word for telling, explaining and describing at B1.
\end{grammarlens}

\subsection*{Your turn}
\begin{itemize}
\raggedright
  \item \emph{Say it:} \emph{obēsus, obēsa, obēsum}
  \item \emph{Say it:} \emph{obēsus, obēsa, obēsum}, once more
  \item \emph{Say it:} a sentence of your own with \emph{obēsus, obēsa, obēsum}
  \item \emph{Recall:} say the Latin for leanness, then the Latin for sturdy, then say \emph{obēsus, obēsa, obēsum} again
\end{itemize}

\begin{tcolorbox}[breakable,colback=teal!4,colframe=teal!35!black,title={Before you move on}]
What does obēsus, obēsa, obēsum mean? (\textquotedblleft{}Fat, stout\textquotedblright{}.) Say it once more.
\end{tcolorbox}
\section[gracilis, gracile]{gracilis, gracile — slender, slim}

\label{lesson:LA-C523-gracilis}



\begin{quote}
Before the new one: say the Latin for sturdy, then the Latin for fat.
\end{quote}

\subsection*{You'll want to know: gracilis, gracile}
\textbf{gracilis, gracile} — \textquotedblleft{}slender, slim\textquotedblright{}.

\begin{quote}
Puella gracilis et alta est.
\end{quote}

\begin{quote}
\emph{The girl is slender and tall.}
\end{quote}

Two endings. Slim in a graceful way; \emph{macer} is thin from want.

\begin{cousinweb}
\textbf{gracilis} is a native Latin adjective; Lewis \& Short compares an old Latin word for \textquotedblleft{}slender\textquotedblright{}, but the origin is not settled. English \emph{gracile} comes from it; \emph{grace} does not.
\end{cousinweb}

\begin{grammarlens}[title={What you've built}]
One more word for telling, explaining and describing at B1.
\end{grammarlens}

\subsection*{Your turn}
\begin{itemize}
\raggedright
  \item \emph{Say it:} \emph{gracilis, gracile}
  \item \emph{Say it:} \emph{gracilis, gracile}, once more
  \item \emph{Say it:} a sentence of your own with \emph{gracilis, gracile}
  \item \emph{Recall:} say the Latin for sturdy, then the Latin for fat, then say \emph{gracilis, gracile} again
\end{itemize}

\begin{tcolorbox}[breakable,colback=teal!4,colframe=teal!35!black,title={Before you move on}]
What does gracilis, gracile mean? (\textquotedblleft{}Slender, slim\textquotedblright{}.) Say it once more.
\end{tcolorbox}
\section[sānitās, sānitātis]{sānitās, sānitātis — health, soundness}

\label{lesson:LA-C523-sanitas}



\begin{quote}
Before the new one: say the Latin for fat, then the Latin for slender.
\end{quote}

\subsection*{You'll want to know: sānitās, sānitātis}
\textbf{sānitās, sānitātis} — \textquotedblleft{}health, soundness\textquotedblright{}.

\begin{quote}
Sānitās maius bonum est quam pecūnia.
\end{quote}

\begin{quote}
\emph{Health is a greater good than money.}
\end{quote}

Feminine, third declension, from \emph{sānus}. Of the mind, good sense. English \emph{sanity}.

\begin{cousinweb}
\textbf{sānitās} is built on \emph{sānus} \textquotedblleft{}sound, healthy\textquotedblright{} with \emph{-itās} (Lewis \& Short; EtymDB). English \emph{sanity} comes from it.
\end{cousinweb}

\begin{grammarlens}[title={What you've built}]
One more word for telling, explaining and describing at B1.
\end{grammarlens}

\subsection*{Your turn}
\begin{itemize}
\raggedright
  \item \emph{Say it:} \emph{sānitās, sānitātis}
  \item \emph{Say it:} \emph{sānitās, sānitātis}, once more
  \item \emph{Say it:} a sentence of your own with \emph{sānitās, sānitātis}
  \item \emph{Recall:} say the Latin for fat, then the Latin for slender, then say \emph{sānitās, sānitātis} again
\end{itemize}

\begin{tcolorbox}[breakable,colback=teal!4,colframe=teal!35!black,title={Before you move on}]
What does sānitās, sānitātis mean? (\textquotedblleft{}Health, soundness\textquotedblright{}.) Say it once more.
\end{tcolorbox}
\section[salūbris, salūbre]{salūbris, salūbre — wholesome, healthy}

\label{lesson:LA-C523-salubris}



\begin{quote}
Before the new one: say the Latin for slender, then the Latin for health.
\end{quote}

\subsection*{You'll want to know: salūbris, salūbre}
\textbf{salūbris, salūbre} — \textquotedblleft{}wholesome, healthy\textquotedblright{}.

\begin{quote}
Aqua huius fontis salūbris est.
\end{quote}

\begin{quote}
\emph{The water of this spring is wholesome.}
\end{quote}

Two endings; an older masculine \emph{salūber} also appears. Of places, food and air. English \emph{salubrious}.

\begin{cousinweb}
\textbf{salūbris} is built on \emph{salūs} \textquotedblleft{}health, safety\textquotedblright{} (Lewis \& Short). English \emph{salubrious} comes from it.
\end{cousinweb}

\begin{grammarlens}[title={What you've built}]
One more word for telling, explaining and describing at B1.
\end{grammarlens}

\subsection*{Your turn}
\begin{itemize}
\raggedright
  \item \emph{Say it:} \emph{salūbris, salūbre}
  \item \emph{Say it:} \emph{salūbris, salūbre}, once more
  \item \emph{Say it:} a sentence of your own with \emph{salūbris, salūbre}
  \item \emph{Recall:} say the Latin for slender, then the Latin for health, then say \emph{salūbris, salūbre} again
\end{itemize}

\begin{tcolorbox}[breakable,colback=teal!4,colframe=teal!35!black,title={Before you move on}]
What does salūbris, salūbre mean? (\textquotedblleft{}Wholesome, healthy\textquotedblright{}.) Say it once more.
\end{tcolorbox}
